% !TEX TS-program = pdflatexmk
\documentclass[a4paper,12pt]{article} %[twocolumn]
% page formatting %%%%%%%%%%%%%%%%%%%%%
%\usepackage{geometry} %defaults are roughly [top=33mm, bottom=48mm, left=33mm, right=33mm]
%\usepackage[top=20mm, bottom=30mm, left=15mm, right=15mm]{geometry}
\usepackage[top=30mm, bottom=40mm, left=25mm, right=25mm]{geometry}
%\usepackage[top=30mm, bottom=45mm, left=30mm, right=30mm]{geometry}
%\usepackage[top=26mm, bottom=40mm, left=26mm, right=26mm]{geometry}

\usepackage{/Users/wdlinch3/Dropbox/Rashoumon/LaTeX/macros/WDL3}

%
%%%
%%%%%
%%%%%%%%%
%%%%%%%%%%%%%%%%%
%%%%%%%%%%%%%%%%%%%%%%%%%%%%%%%%%
%%%%%%%%%%%%%%%%%%%%%%%%%%%%%%%%%%%%%%%%%%%%%%%%%%%%%%%%%%%%%%%%%
%
%UMDEPP-014-026
%\title{Components of Eleven-dimensional Supergravity from the Superspace Chern-Simons Action}
%%%%%%%%%%%%%%%%%%%%%%%%%%%%%%%%%%%%%%%%%%%%%%%%%%%%%%%%%%%%%%%%
\begin{document}
\nopagecolor
%\maketitle

%%\begin{center}
%\addtolength{\baselineskip}{.5mm}
%\thispagestyle{empty}
%\begin{flushright}
%{\color{red}\today \\
%{\sc MI-TH-16??}\\
%}
%\end{flushright}

\vspace{10mm}
\begin{center}
%\doublespacing{
\Large \textbf {
Off-shell Supersymmetry and the M-theory Effective Action
}
%}
\end{center} 
\vspace{5mm}
\begin{center} 
%\\[10mm] 
%\\[15mm]
{
William D.\ Linch III
}
\\[5mm]
%~\\[5mm]
{
{\em George P. and Cynthia W. Mitchell Institute for\\
Fundamental Physics and Astronomy, \\
Texas A\&{}M University.}\\
~\\
11 September, 2019.\\
~\\
Based on work 
%\cite{Becker:2016xgv, %ATH
%Becker:2016rku,	%NATH
%Becker:2016edk,	%G2Potential
%Becker:2017njd, 	%All CS
%Becker:2017zwe,	%Linearized
%Becker:2018phr}	%Current
with \\
~\\
K.\ Becker, 
M.\ Becker, 
D.\ Butter, 
S.\ Guha, 
S.\ Randall, 
D.\ Robbins,
and
A.\ Sengupta.
}
\end{center} 

\vspace{10pt}


\begin{abstract}
These are notes for a seminar in the Theory Group at Masaryk University delivered on 8 October, 2019.
\end{abstract}

%%email addresses
%%\vspace*{.5cm}
%\begin{flushleft}
%~\\
%{${}^{\small \mbox\Pisces}$ \href{mailto:wdlinch3@gmail.com}{wdlinch3@gmail.com}}\\
%%\makebox[0pt][t]
%%{$^\text{\Scorpio}$ \href{mailto:siegel@insti.physics.sunysb.edu}{siegel@insti.physics.sunysb.edu}}
%%\makebox[0pt][t]
%\end{flushleft}
\thispagestyle{empty}
\newpage
%\thispagestyle{empty}
%\tableofcontents

%\newpage
\setcounter{page}1
%
%%%%%%%%%%%%%%%%%%%%%%%%%%%%%%%%%%%%%%%%%%%%%%%%%%%%%%%%%%%%%%%%
%%%%%%%%%%%%%%%%%%%%%%%%%%%%%%%%%%%%%%%%%%%%%%%%%%%%%%%%%%%%%%%%
\section{Motivation and Strategy}
Quantum field theory may be described accurately as a calculus for relativistic quantum fields (spins $\leq 1$).\\
(The theory formerly known as) String Theory forms the rudiments of a calculus for quantum gravity (spin 2) coupled to quantum fields (spins $<2$). \\
Currently, its simplest formulation is in terms of first quantized free superstrings in ten dimensions.\\
In this perturbative formulation, these are not unique. Non-perturbatively, they are not complete. \\
In 1995 Witten conjectures that there is a unique non-perturbative completion with an eleven-dimensional supergravity limit in its moduli space.
\begin{displaymath}
\xymatrix@=10pt{ 
&\ar@{-}@/^1pc/[dl]
	{\begin{array}{c}11\textrm{\sc d}\, \textrm{{\sc sg}}\end{array}} 
	\ar@{-}@/_1pc/[dr]& \\
\mathbf {II}\textrm{\sc a}\ar@{-}@/^1pc/[dd]&&\mathbf {H}\textrm{\sc e}\ar@{-}@/_1pc/[dd]\\
&\mathbf M&\\
\mathbf {II}\textrm{\sc b}&&\mathbf {H}\textrm{\sc o}\\
&\ar@{-}@/_1pc/[ul]{\mathbf I}\ar@{-}@/^1pc/[ur]&
}
\end{displaymath}
This ``M-theory'' is still poorly understood.
It should have (at least) 2-branes and 5-branes (not-necessarily-fundamental). These have comparable tensions (not parametrically-separated), so a world-volume-like description has not been found.\\
The alternative is to exploit the spacetime properties (\eg symmetries and anomaly structure) together with known limits in the moduli space of theories (\ie string theories). 

%%%%%%%%%%%%%%%%%%%%%%%%%%%%%%%%%%%%%%%%%%%%%%%%%%%%%%%%%%%%%%%%
\subsection{Noether Procedure}
We will focus on the following three properties of M:
\begin{enumerate}
\item 11D super-Poincar\'e symmetry represented by the minimal multiplet as explained by Cremmer, Julia and Scherk (1978): 11D SG represented in components by a metric $g_{\bm {mn}}$, a gravitino $\psi_{\bm m}^{\bm \alpha}$, and a gauge 3-form $C_{\bm {mnp}}$. 
\item Partial knowledge of the leading higher-derivative corrections to the 11D SG: $\int C \wedge X_8$, with $X_8 \sim R^4+\cdots$ a specific contraction of the fourth exterior power of the curvature 2-form with corrections involving fewer curvatures.
\item The 4-form field strength $G=dC$ does not have integer periods on a generic spinnable Lorenzian 11-manifold. What has integral periods is instead $G - \tfrac14 p_1$ where $p_1 \sim \mathrm{tr} \, R\wedge R$ is the first Pontryagin class (Witten 1997).
\end{enumerate}
We may represent the situation as follows: The M-theory effective action is of the form $S=S_2 + S_8 +\cdots$ where the subscript denotes the number of derivatives on the bosonic fields. \\
Then under the lowest-order supersymmetry transformation $\delta_0$, $S_2$ is invariant, but $S_8$ is not. \\
Supersymmetry of the effective action requires a correction $\delta= \delta_0+\delta_6 +\cdots$ so that $\delta_6 S_2$ cancels the transformation $\delta_0 S_8$. \\
But then $\delta_6 S_8\neq 0$ so that we require an even-higher-order term in the effective action that may cancel this, and so on.
This is the usual component Noether procedure--it has been carried out explicitly only to very low orders and even then only partially.\\
Putting this enormous technical hurdle aside, summing up this infinite number of steps would give a supersymmetric, all-orders-in-derivatives effective action. \\
It is a long-standing conjecture that this is the leading Wilsonian effective action of M-theory. Schematically,
\begin{align}
\label{E:MLEEA}
R + R^4 + \textrm{11D, $N=1$ supersymmetry} \Rightarrow \textrm{Leading M-theory LEEA?}
\tag{$\ast$}
%\nonumber
\end{align}

%%%%%%%%%%%%%%%%%%%%%%%%%%%%%%%%
\paragraph{Born-Infeld}

The situation described above has a spin-1 analog discovered by Born and Infeld (1933). 
They deformed the Maxwell Lagrangian by higher-derivative corrections in such a way that the field strength would be bounded above: $L\sim \frac1{\alpha^2 }[1-\sqrt{1-\alpha^2 F^2}]\approx \tfrac12 F^2 + o(\alpha^2 )$.\\
It was discovered by Tseytlin {\em half a century} later (1985) that this action is the exact leading LEEA of open (super)string field theory. \\
A decade later, it was realized by Bagger and Galperin (1997) that this action is fixed by extended supersymmetry, so that we have the analog 
\begin{align}
F^2 + F^4 + \textrm{10D, $N=1$ supersymmetry} \Rightarrow \textrm{Born-Infeld}
\nonumber
\end{align}
of the conjecture \eqref{E:MLEEA} above. 



%%%%%%%%%%%%%%%%%%%%%%%%%%%%%%%%%%%%%%%%%%%%%%%%%%%%%%%%%%%%%%%%
\subsection{Off-shell Supersymmetry}
The Noether-type situation described above is typical of theories with ``soft'' supersymmetry algebras: $\{ Q, Q\}\sim P$ on the fields only up to the lower-order equations of motion, or ``on-shell''.\\
Actual algebras close on field representations irrespective of this; for contrast, they are referred to as ``off-shell''. \\
In such representations, there are additional ``auxiliary'' component fields that do not propagate, but will be related to interactions terms through their equations of motion. \\
When the off-shell formulation exists, the auxiliary fields can (in principle) be integrated out which changes the interaction terms in the Lagrangian and supersymmetry transformations thereby reducing to the off-shell theory to the on-shell formulation. 

Now suppose there were an off-shell, 11D superspace with all 32 supercharges. \\
Then a description of the leading terms in the M-theory LEEA in this superspace must actually contain all orders, at least when all the auxiliary fields are integrated out. \\
The catch is that this superspace does not exist. (For rigid supersymmetry Berkovits (1993) explains that the algebra will not close unless the number of supercharges is 9 or less due to non-associativity of $\mathbb O$.)

Can we make do with fewer than 32 supercharges, say, only 4?\\
In this 11D, $N=1/8$ superspace, the 4 supersymmetries would close on 4 translations giving a 4D, $N=1$ {\em off-shell} subalgebra of the full 11D, $N=1$ on-shell supersymmetry algebra.\\
We would then have to impose the remaining symmetries ``by hand''. \\
Let us make this more concrete with a very good analogy: 

%%%%%%%%%%%%%%%%%%%%%%%%%%%%%%%%%%%%%%%%%%%%%%%%%%%%%%%%%%%%%%%%
%%%%%%%%%%%%%%%%%%%%%%%%%%%%%%%%%%%%%%%%%%%%%%%%%%%%%%%%%%%%%%%%
\section{Super-Maxwell in 10D, $N=1/4$ Superspace}

Consider 10D super-Maxwell theory but, again, with only 4 supercharges manifest, that is, rigid 10D, $N=1/4$ superspace. 
The Maxwell multiplet $(A_{\bm m}, \lambda^{\bm \alpha})$ breaks up under the $\mathrm D=4+6$ split into
\begin{align}
(A_m , \lambda^\alpha) 
~,~~ 
(\varphi_i , \psi^\alpha_i)
~,~\and~~ 
(\bar \varphi^i , \bar \psi_{\dt \alpha}^i)
\nonumber
\end{align}
with indices $m=0,1,2,3$ for 4D coordinate, $\alpha=1,2$ for 4D Weyl, and $i = 1,2,3$ for the defining representation of $SU(3)$. 
We may embed these into 4D, $N=1$ superfields in the usual way
\begin{align}
V(x, \theta, \bar \theta;y) & \sim \cdots + (\theta \sigma^m \bar \theta) A_m 
	+ \bar \theta^2 \theta^\alpha \lambda_\alpha 
	+ \theta^2 \bar \theta_{\dt \alpha} \bar \lambda^{\dt \alpha}
	+\theta^2 \bar \theta^2 D
\cr
\Phi_i(x, \theta;y) & \sim \varphi_i + \theta^\alpha \psi_{\alpha i} + \theta^2 F_i
~.
\nonumber
\end{align}
Under the 10D, $N=1$ gauge transformations,
\begin{align}
\delta V = \frac1{2i} (\Lambda -\bar \Lambda)
~~~\and~~~
\delta \Phi_i = \partial_i \Lambda
~.
\nonumber
\end{align}
The invariant field strengths are easily checked to be
\begin{align}
W^\alpha &= -\frac14 \bar D^2 D^\alpha V
\cr
Z^\alpha_i &= D^\alpha(\Phi_i - 2i\partial_i V)
\cr 
F_i{}^j &= (\partial_i \bar \Phi^j - \bar \partial^j \Phi_i)  + 2i \bar \partial^j \partial_i V
\cr
E_{ij} &= \partial_i \Phi_j -\partial_j \Phi_i
\nonumber
\end{align}
The Lagrangian can no longer be written in manifestly supersymmetric form:
\begin{align}
L =  -\frac 18 \int d^2\theta \,\Phi_i \left[
	 \bar{D}_{\dt \alpha} \bar {Z}^{\dt \alpha i} 
	-{2\sqrt{2}} \varepsilon^{ijk} E_{jk} 
	 \right]
	-\frac14  \int d^4\theta \,V \left[
	 D^\alpha W_\alpha 
	-2i F_k{}^k
	\right]
	+\hc
\nonumber
\end{align} 
This is a Chern-Simons-type Lagrangian with only the third term writeable in manifestly gauge-invariant form.\\
In components,
\begin{align}
\mathcal L_{max} &= 
-\tfrac14 F^{ab}F_{ab} + \tfrac12 d^2
-F^{a k} F_{a k} + \bar f^k f_k - i d F_k{}^k
-\tfrac i{\sqrt2}\varepsilon^{ijk} f_i E_{jk} 
+\tfrac i{\sqrt2}\varepsilon_{ijk} \bar f^i\bar E_{jk} 
	+\mathrm{fermions}
~.
\nonumber
\end{align}
The auxiliary field equations of motion are 
\begin{align}
%\label{E:AuxFieldsEoM1}
d \approx i  F_k{}^k
~~~\and~~~
\bar f^k \approx  \frac{i}{\sqrt2} \varepsilon^{ijk} E_{ij}
~.
\nonumber
\end{align}
Substituting these algebraic equations back into the component action we get
\begin{align}
\mathcal L_{max} &= -\frac14 F^{ab}F_{ab} 
	-  F^{a k} F_{a k}
	+\frac12 (F_k{}^k)^2
	-\bar E^{ij} E_{ij} 
	+\mathrm{fermions}
~.
\nonumber
\end{align}
This is not the 10D Maxwell Lagrangian, but integrating twice by parts, 
\begin{align}
\int d^6 y \, (F_k{}^k)^2 = \int d^6y \, [ F_i{}^j F_j{}^i +\bar E^{ij} E_{ij} ]
\nonumber
\end{align}
all the pieces reassemble into ten-dimensional super-Maxwell:
\begin{align}
\mathcal L_{max} &= -\frac14 F^{ab}F_{ab} 
	-  F^{a k} F_{a k}
	+\frac12 F_i{}^j F_j{}^i
	-\frac12\bar E^{ij} E_{ij} 
	+\mathrm{fermions}
~.
\nonumber
\end{align}

%%%%%%%%%%%%%%%%%%%%%%%%%%%%%%%%%%%%%%%%%%%%%%%%%%%%%%%%%%%%%%%%
\subsection{Toward Born-Infeld}

Now suppose we would like to understand the Born-Infeld action in this superspace. 
Firstly, the action in 4D, $N=1$ is known: It is given by Cecotti and Ferrara (1987) as
%\begin{align}
%L_{BI} = \frac 14 \int d^2\theta \, \bm X
%	+\hc
%~\st~
%\bm X=  W^2 +\frac12 \bar D^2 \left[ { W^2 \bar W^2 \over 1 +2 \hat T + \sqrt{1 +4(\hat T -\check T^2)}}\right]
%\nonumber
%\end{align} 
\begin{align}
\mathcal L_{CF} = \frac 12 \int d^2\theta \, W^2
	- \int d^4\theta \, { W^2 \bar W^2 \over 1 +2 \hat T + \sqrt{1 +4\hat T - 4\check T^2}}
%~\st~
%\bm X=  W^2 +\frac12 \bar D^2 \left[ { W^2 \bar W^2 \over 1 +2 \hat T + \sqrt{1 +4(\hat T -\check T^2)}}\right]
\nonumber
\end{align} 
where $W^2 =W^\alpha W_\alpha$ is the super-photon field strength superfield and  $T:= -\frac14 \bar D^2 (\bar W^2) = d^2 -\frac12 F_{ab} [ F_{ab} +i \tilde F^{ab}  ]+$ fermions.\\
This first term is the 4D Maxwell action; we lift it to 10D as above. (Note that this is not just a replacement of field strengths.) \\
The leading correction in the field strength expansion is $W^2 \bar W^2$. This can be lifted to
\begin{align}
W^2\bar W^2 \longrightarrow W^2\bar W^2 + \alpha (WZ_i)(\bar W\bar Z^i) + \beta Z^2_{ij} \bar Z^{2\, ij}
\nonumber
\end{align} 
for some $\alpha, \beta \in \mathbf R$.
This with reproduce the correct mixed component terms (after integrating out auxiliary fields) provided $\alpha = -4$ and $\beta = 1$, but it does not give the ``all-internal'' parts. This is a loose end.

Focusing further on the mixed terms, we can contemplate non-linear corrections by turning off the 4D vector completely: set $V\to 0$. 
This reduces the action to just the scalars. This action should match that of Ro\v{c}ek and Tseytlin (1997):
\begin{align}
\mathcal L_{RT} 
%	= \int d^4\theta \left[
%	\bar \Phi^i \Phi_i + {Z_{ij}^2 \bar Z^{2ij} \over 
%	1+ \frac A2 + \sqrt{1+A+B}
%	}
%	\right]
	= \int d^4\theta
	\bar \Phi^i \Phi_i 
	+\int d^4\theta {Z_{ij}^2 \bar Z^{2ij} \over 
	1+ \frac A2 + \sqrt{1+A+B}
	}
\nonumber
\end{align}
with 
\begin{align}
A= -4 [ F^a{}_i \bar F_a{}^i  + f_i \bar f^i ] 
~~~\and~~~
%B= 4 (F^a{}_i \bar F_a{}^i)^2 - 4 (F^a{}_i F_{a j})(\bar F_b{}^i \bar F^{bj})
B= 4 (F^a{}_i \bar F_a{}^i)(F^b{}_j \bar F_b{}^j) - 4 (F^a{}_i F_{a j})(\bar F_b{}^i \bar F^{bj})
\nonumber
\end{align}
This suggests that we further extend 
\begin{align}
\mathcal L_{max} + \int d^4\theta \, {W^2\bar W^2 -16 (WZ_i)(\bar W\bar Z^i) +4 Z^2_{ij} \bar Z^{2\, ij} \over 
	1 +2 \hat {\bm T} + \sqrt{1 +4(\hat {\bm T} -\check {\bm T}^2)}
}
\end{align}
with ${\bm T} = -\frac14 \bar D^2 \bar W^2 +\tfrac14 A -\frac i4 B +...$
%\begin{align}
%{\bm T} &= -\frac14 \bar D^2 \bar W^2 +\tfrac14 A -\frac i4 B +...
%\cr
%%	&= -\frac14F^{\s+2}_{ab} +\frac12 d^2 
%%	-F^a{}_i \bar F_a{}^i + f_i \bar f^i  - i d F_i{}^i  -\frac i{\sqrt2}\varepsilon^{ijk} f_i E_{jk} 
%%		+ ...
%	&= -\frac14F^{\s+2}_{ab} -F^a{}_i \bar F_a{}^i 
%	+\frac12 d\left(d-2i  F_i{}^i\right) + f_i \left(\bar f^i   -\frac i{\sqrt2}\varepsilon^{ijk}  E_{jk} \right)
%		+ ...
%\end{align}

This is unfinished work. Instead of exploring it further here, we push it over the cliff... 

%%%%%%%%%%%%%%%%%%%%%%%%%%%%%%%%%%%%%%%%%%%%%%%%%%%%%%%%%%%%%%%%
%%%%%%%%%%%%%%%%%%%%%%%%%%%%%%%%%%%%%%%%%%%%%%%%%%%%%%%%%%%%%%%%
\section{Putting it up to 11(D SG in $N=1/8$ Superspace)}
Following the strategy for spin-1, we first decompose 11D SG fields 
\begin{align}
\begin{array}{cccclcc}
	e_{\bm m}{}^{\bm a} &\to &
		e_{m}{}^{a} ~,~ 
		\mathcal A_m^i ~,~ 
		{\color {blue} g_{ij} }
&\with & a, m = 0,1,2,3
\\
	\psi_{\bm m}^{\bm \alpha}&\to &
		\psi_{m}^{\alpha} ~,~ 
		\psi_{m}^{\alpha i} ~,~ 
		\psi_{j}^{\alpha } ~,~ 
		\psi_{j}^{\alpha i} ~,~
&\and & i,j = 1,\dots, 7	
\\ 
	C_{\bm {mnp}} &\to &C_{mnp}~,~ C_{mn \, i} ~,~ C_{m \, ij}~,~ C_{ijk}
&\and & \alpha = 1,2	
\end{array}
\nonumber
\end{align}
and embed them into superfields:
\begin{align}
U^a &\sim \cdots 
	+ (\theta \sigma^m \bar \theta) e_m{}^a 
	+ \theta^2 (\sigma^m \bar \theta)_\alpha \psi_m{}^\alpha 
	+ \bar \theta^2 (\theta \sigma^m )_{\dt \alpha} \bar \psi_m{}^{\dt \alpha }
	+ \cdots 
\cr
\Psi^{\alpha i} &\sim \cdots 
	+ (\theta \sigma^m \bar \theta) \psi_m{}^{\alpha i}
	+\cdots
\cr
\mathcal V^i &\sim \cdots 
	+ (\theta \sigma^m \bar \theta) \mathcal A_m^i 
%	+ \bar \theta^2 \theta^\alpha \lambda_\alpha^i
%	+ \theta^2 \bar \theta^{\dt \alpha}\bar \lambda_{\dt \alpha}^i
	+\cdots
	+ \theta^2\bar \theta^2 \bm d^i
\cr
X &\sim \cdots 
	+ \theta^2 ({\color{red}s}+i{\color{red}p}) 
	+ \bar \theta^2 ({\color{red}s}-i{\color{red}p}) 
	+ \theta \sigma_q \bar \theta \mathcal \, \varepsilon^{mnpq} C_{mnp}
	+ \cdots
	+ \theta^2\bar \theta^2 d_X
\cr
\Sigma_i^\alpha &\sim \cdots	
	+ \theta^\beta [\delta_\beta^\alpha {\color{red}h_i} + (\sigma^{mn})_\beta{}^\alpha C_{mn \, i} ] 
	+  \cdots
\cr
V_{ij} &\sim \cdots 
	+  (\theta \sigma^m \bar \theta)  C_{m\, ij}
%	+ \bar \theta^2 \theta^\alpha \lambda_\alpha^i
%	+ \theta^2 \bar \theta^{\dt \alpha}\bar \lambda_{\dt \alpha}^i
	+\cdots
	+ \theta^2\bar \theta^2 d_{ij}
\cr
\Phi_{ijk} &\sim C_{ijk} + i {\color{blue} \varphi_{ijk} } + \cdots + \theta^2 f_{ijk}
\nonumber
\end{align}
This requires the following surprising fact about Riemannian 7-manifolds: \\
For any orientable and spinnable such smooth manifold, 
%\begin{theorem}
%Let $(Y, g)$ be a smooth Riemannian manifold that is oriented ($w_1=0$) and spin ($w_2$=0). 
%Then 
%\end{theorem}
there is a 3-form $\varphi$ such that 
\begin{align}
\sqrt{g}g_{ij} = \frac1{144}\epsilon^{k_1\cdots k_7} \varphi_{i k_1k_2} \varphi_{j k_3 k_4} \varphi_{k_5k_6k_7}
\nonumber 
\end{align}

%%%%%%%%%%%%%%%%%%%%%%%%%%%%%%%%%%%%%%%%%%%%%%%%%%%%%%%%%%%%%%%%
\subsection{Chern-Simons-like Action}
The next step is to supersymmetrizing the term $\int C \wedge dC \wedge dC$ in the M-theory effective action. \\
The spins $\leq 1$ fields in the list form a ``non-abelian tensor hierarchy'' of superfields.\footnote{Mathematically, this is a {\it Diff}(7)-equivariant chain complex of differential forms on $Y$ with values in a complex of (sheaves of) representations of the 4D, $N=1$ super-Poincar\'e algebra.} 
Such homological objects have a complex of field strengths:
\begin{align}
E_{ijkl} &= 4\partial_{[i} \Phi_{jkl]}
	\cr
F_{ijk} &= \frac1{2i}\left( \Phi_{ijk} - \bar \Phi_{ijk}\right) - 3\partial_{[i} V_{jk]}
	\cr
W_{\alpha ij} &= -\frac14 \bar \nabla^2 \nabla_\alpha V_{ij}
	+2\partial_{[i} \Sigma_{\alpha j]}
	+ \mathcal W_\alpha^k \Phi_{ijk}
	\cr
H_i &= \frac1{2i}\left(\nabla^\alpha \Sigma_{\alpha i} - \bar \nabla_{\dt \alpha} \bar \Sigma_i^{\dt \alpha} \right)
	-\partial_i X 
	+\left[  \mathcal W^{\alpha j} \nabla_\alpha  V_{ij}
		+\tfrac12 \nabla^\alpha \mathcal W_\alpha^j   V_{ij}
	+\hc \right]
	\cr
G&= -\frac14 \bar \nabla^2 X 
	+ \mathcal W^{\alpha i} \Sigma_{\alpha	i}
\nonumber
\end{align}
They also have a unique Chern-Simons-like invariant:
%\begin{align}
%	-12 \kappa^2 \mathcal L_{CS} &= i\int d^2 \theta \, \left[ i \Phi_{[0,3]}\wedge  \mathbb G_{[4,4]}  + \Sigma^\alpha_{[2,1]} \wedge \mathbb W_{[2,6] \alpha }\right] 
%	\cr&
%	+\int d^4 \theta  \left[ V_{[1,2]} \wedge \mathbb H_{[3, 5]} - X_{[3,0]} \mathbb F_{[1,7]}\right]
%	+\hc 
%\nonumber
%\end{align}
\begin{align}
-12 L_{CS} &= 
i \int d^2 \theta  \,   \Phi \left[
	E G + \tfrac12 W^\alpha W_\alpha - \tfrac i4 \bar {\nabla}^2 (  F H) 
	\right]
\cr&+ \int d^4 \theta  \,  V \left[  
	\left( E + \bar E\right) H + \omega(W, F) 
	- i {\nabla}^\alpha F\iota_{\mathcal W_\alpha}F 
	+ i \bar {\nabla}_{\dt \alpha} F \iota_{\bar{\mathcal W}^{\dt \alpha} }F 
	\right]
\cr&+i \int d^2 \theta  \, \Sigma^\alpha \left[
	E W_\alpha -\tfrac i4 \bar {\nabla}^2 (  F {\nabla}_\alpha F) 
	\right]
\cr&- \int d^4 \theta  \, X \left[  
	\left( E + \bar E\right) F
	\right]
	+\mathrm{h.c.} 	
\nonumber
\end{align}
Some observations:
\begin{itemize}
	\item 
		$\Phi\wedge E\wedge G = G \Phi \wedge d \Phi $ contains the ``$G_2$ superpotential'' (Gukov 1999). 
	\item 
		$\Phi\wedge \mathbb G = \frac12 \Phi \wedge W^\alpha \wedge W_\alpha +\cdots$ also contains the holomorphic gauge-kinetic function $f(\Phi)  = \Phi$ familiar from Seiberg-Witten theory.
	\item Every term is required by gauge invariance: Had we not known this action existed, we would have discovered it by starting from any single particular term.
	\item Local 4D, N=1 super-conformal invariance.
\end{itemize}


%%%%%%%%%%%%%%%%%%%%%%%%%%%%%%%%%%%%%%%%%%%%%%%%%%%%%%%%%%%%%%%%
\subsection{Supervolume}
Local supersymmetry on 4D, $N=1$, diffeomorphism invariance in 7D, tensor hierarchy gauge invariance, engineering dimensions, conformal weights, $U(1)_R$ weights, and $GL(7)$ representation define the volume of 11D, $N=1/8$ superspace 
\begin{align}
	S_H = -\frac 3{\kappa^2} \int d^7 y \int d^4x  \int d^4\theta \, \mathrm{sdet}(E_M{}^A) (\bar GG)^{1/3} \sqrt{g(F)}\, f(x)
%	~~~,~~~ x = {H_i g^{ij}(F) H_j \over (\bar GG)^{2/3}}
\nonumber
\end{align} 
up to a function of $x = {H_i g^{ij}(F) H_j \over (\bar GG)^{2/3}}$.\\
Introducing the gravitino superfield $\Psi^\alpha_i$, we can extend the local supersymmetry from 4D, $N=1$ to all of 11D, $N=1/8$. This fixes $f(x)$ exactly: 
The combination ${\hat f} := f -2 x f'$ satisfies the quartic equation 
\begin{align}
	\frac x4 {\hat f}^4 +  {\hat f}^3 -1 = 0
		~~~\and~~~
	f(0)= 1
	~~~\Rightarrow
	~~~f(x) = \textrm{ugly but explicit function.}
\nonumber
\end{align}


%%%%%%%%%%%%%%%%%%%%%%%%%%%%%%%%%%%%%%%%%%%%%%%%%%%%%%%%%%%%%%%%
\subsection{Components}
A big mystery is how in the world this can produce the correct ``scalar potential'' 
\begin{align}
\label{E:ScalarPotential}
V(x,y)&= 
	-\frac12 R(g)
		+\frac1{4\cdot 4!} F_{ijkl}^2 
\tag{$\ast \ast$}
\end{align}
$=$ Einstein-Hilbert and Maxwell terms with all polarizations along 7D.\\
As with Maxwell, this must happen by (more complicated) auxiliary field mechanism: 
In terms of these fields, the scalar potential is 
\begin{align}
%\label{E:ScalarPotential1}
\hspace{-8mm}V=
	\tfrac14d_X^2 
	+ 2 g_{ij} \bm d^i \bm d^j 		
	-  d_{\mathbf 7 ij} ^2 
		+ \tfrac12 d_{\mathbf {14} ij}^2
	-\tfrac1{18} |f_{\mathbf 1 ijk}+\tfrac{i}2d_XF_{ijk}|^2
		-\tfrac1{24} |f_{\mathbf 7 ijk}|^2
		+\tfrac1{24} |f_{\mathbf {27} ijk}|^2
\nonumber
\end{align}
Define the ``torsion forms'' of $\varphi$ by 
\begin{align}
d\varphi = \tau_0 \ast \varphi + 3\tau_1 \wedge \varphi + \ast \tau_3
~~~\and~~~
d(\ast \varphi) = 4\tau_1 \wedge \ast \varphi + \tau_2 \wedge \varphi
\nonumber
\end{align}
and define $\sigma_0,\cdots, \sigma_3$ similarly for the 3-form $C_{ijk}$.\\
Then component projection results in the equations of motion for the auxiliary fields. Plugging them in gives
\begin{align}
%\label{E:ScalarPotential}
V(x,y)&= 
	-\frac{21}{16} \tau_0^2 
	- 15\tau_1^2 
	+\frac18 \tau_2^2 
	+\frac1{24} \tau_3^2			
%
	+\frac{7}{4} \sigma_0^2 
	+ 9 \sigma_1^2  
	+\frac1{24} \sigma_3^2	
=\eqref{E:ScalarPotential}	
\nonumber
\end{align}

%%%%%%%%%%%%%%%%%%%%%%%%%%%%%%%%%%%%%%%%%%%%%%%%%%%%%%%%%%%%%%%%
\subsection{Beyond...}




%%%%%%%%%%%%%%%%%%%%%%%%%%%%%%%%%%%%%%%%%%%%%%%%%%%%%%%%%%%%%%%%
%%%%%%%%%%%%%%%%%%%%%%%%%%%%%%%%%%%%%%%%%%%%%%%%%%%%%%%%%%%%%%%%
%%%%%%%%%%%%%%%%%%%%%%%%%%%%%%%%%%%%%%%%%%%%%%%%%%%%%%%%%%%%%%%%
%%%%%%%%%%%%%%%%%%%%%%%%%%%%%%%%%%%%%%%%%%%%%%%%%%%%%%%%%%%%%%%%
%{
%%\small
%\footnotesize
%%\scriptsize
%%
%\bibliography{/Users/wdlinch3/Dropbox/Rashoumon/LaTeX/BibTex/BibTex}
%\bibliographystyle{unsrt}
%%\bibliographystyle{amsplain}
%%\bibliographystyle{amsalpha}
%}

\end{document}  


